\corrigesChapitre{19}{Le son : fréquence, propagation et écho}
\par\noindent\begin{minipage}{\linewidth}
\begin{corrige}{1}{Le milieu}
\begin{questions}
\item Non, il nécessite un milieu matériel.
\item Oui, avec une vitesse qui diffère de celle dans l’air.
\end{questions}
\end{corrige}
\end{minipage}
\par\noindent\begin{minipage}{\linewidth}
\begin{corrige}{2}{La source}
\begin{questions}
\item Une vibration transmise au milieu.
\item Non : elles oscillent localement et transmettent la perturbation.
\end{questions}
\end{corrige}
\end{minipage}
\par\noindent\begin{minipage}{\linewidth}
\begin{corrige}{3}{L’unité}
\begin{questions}
\item Le hertz, Hz.
\item Cent motifs ou cycles par seconde.
\end{questions}
\end{corrige}
\end{minipage}
\par\noindent\begin{minipage}{\linewidth}
\begin{corrige}{4}{Trois domaines}
\begin{questions}
\item Infrason, son dans la plage audible humaine usuelle, ultrason.
\item Non : elle concerne la fréquence.
\end{questions}
\end{corrige}
\end{minipage}
\par\noindent\begin{minipage}{\linewidth}
\begin{corrige}{5}{Le plus aigu}
\begin{questions}
\item Celui de 800 Hz.
\item Non : la fréquence ne donne pas à elle seule le niveau sonore.
\end{questions}
\end{corrige}
\end{minipage}
\par\noindent\begin{minipage}{\linewidth}
\begin{corrige}{6}{Une distance simple}
\begin{questions}
\item $d=340\times2=680\unite m$.
\item Non, c’est un trajet simple de la source au récepteur.
\end{questions}
\end{corrige}
\end{minipage}
\par\noindent\begin{minipage}{\linewidth}
\begin{corrige}{7}{Un écho}
\begin{questions}
\item $d=340\times0{,}50/2=85\unite m$.
\item $170\unite m$ au total.
\end{questions}
\par\smallskip\noindent\textbf{Pour comprendre.}\enspace Le produit $340\times0{,}50=170\unite m$ donne le trajet complet. La paroi est à $d=170/2=85\unite m$. Écrire le facteur deux avant d’effectuer le calcul évite de confondre les deux distances.\par
\pcFigureEx{figures/cor-echo-sonore.pdf}{Le son parcourt 170 m au total ; la paroi est à 85 m.}
\end{corrige}
\end{minipage}
\par\noindent\begin{minipage}{\linewidth}
\begin{corrige}{8}{Compter des motifs}
\begin{questions}
\item $f=8/0{,}020=400\unite{Hz}$.
\item Le son simple de 400 Hz est plus aigu, à conditions de comparaison adaptées.
\end{questions}
\par\smallskip\noindent\textbf{Pour comprendre.}\enspace Convertir 20 ms en 0,020 s puis $f=N/\Delta t=8/0{,}020=400\unite{Hz}$. Le nombre de motifs est un comptage, pas une durée. À 400 Hz, le son simple est plus aigu qu’à 200 Hz.\par
\pcFigureEx{figures/cor-huit-motifs.pdf}{On compte huit motifs complets dans 20 ms, soit 0,020 s ; ne pas confondre nombre de crêtes et nombre d’intervalles.}
\end{corrige}
\end{minipage}
\par\noindent\begin{minipage}{\linewidth}
\begin{corrige}{9}{Le tonnerre}
\begin{questions}
\item $d\approx340\times5=1\,700\unite m$.
\item Le temps de propagation de la lumière, très faible devant celui du son sur cette distance.
\end{questions}
\end{corrige}
\end{minipage}
\par\noindent\begin{minipage}{\linewidth}
\begin{corrige}{10}{Dans l’eau}
\begin{questions}
\item $d=1\,500\times0{,}20/2=150\unite m$.
\item Le milieu est l’eau ; la vitesse fournie est adaptée à ce milieu.
\end{questions}
\par\smallskip\noindent\textbf{Pour comprendre.}\enspace Utiliser la vitesse dans l’eau, puis le trajet double : $d=v\Delta t/2=1\,500\times0{,}20/2=150\unite m$. Employer la vitesse dans l’air changerait le modèle physique.\par
\pcFigureEx{figures/cor-sonar.pdf}{La durée est celle d’un aller-retour dans l’eau : la vitesse fournie est 1 500 m/s.}
\end{corrige}
\end{minipage}
\par\noindent\begin{minipage}{\linewidth}
\begin{corrige}{11}{Un signal plus grand}
\begin{questions}
\item Non : la répétition temporelle est identique.
\item La hauteur est identique ; les amplitudes et niveaux peuvent différer selon la mesure et le dispositif.
\end{questions}
\pcFigureEx{figures/cor-meme-periode.pdf}{Même durée d’un motif, donc même fréquence et même hauteur pour ces sons simples ; amplitudes différentes.}
\end{corrige}
\end{minipage}
\par\noindent\begin{minipage}{\linewidth}
\begin{corrige}{12}{Une information incomplète}
\begin{questions}
\item Non : la durée et la répétition comptent aussi, ainsi que les conditions de mesure.
\item Durée d’exposition, répétition, position du récepteur et variations de niveau ; on ne réalise pas de test fort pour vérifier.
\end{questions}
\end{corrige}
\end{minipage}
\par\noindent\begin{minipage}{\linewidth}
\begin{corrige}{13}{Deux trajets confondus}
\begin{questions}
\item Pour l’écho, la distance de la paroi est $102/2=51\unite m$.
\item Oui : sur un trajet simple, $d=v\Delta t=102\unite m$.
\end{questions}
\pcFigureEx{figures/cor-echo-030.pdf}{102 m représente le trajet aller-retour ; la distance à la paroi est 51 m.}
\end{corrige}
\end{minipage}
\par\noindent\begin{minipage}{\linewidth}
\begin{corrige}{14}{Un délai instrumental}
\begin{questions}
\item Délai de propagation : $0{,}200\unite s$ ; $d=340\times0{,}200/2=34\unite m$.
\item Le retard électronique ne correspond pas à un trajet du son et doit être retiré du temps total d’abord.
\end{questions}
\par\smallskip\noindent\textbf{Pour comprendre.}\enspace Durée réelle de propagation : $0{,}210-0{,}010=0{,}200\unite s$. Distance : $d=340\times0{,}200/2=34\unite m$. Le retard électronique n’est pas un temps passé à parcourir l’aller-retour.\par
\pcFigureEx{figures/cor-echo-retard.pdf}{Retirer le retard instrumental d’abord ; seule la durée de propagation corrigée sert à calculer la distance.}
\end{corrige}
\end{minipage}
\par\noindent\begin{minipage}{\linewidth}
\begin{corrige}{15}{Un contrôle de vraisemblance}
\begin{questions}
\item Non : 2 s correspondraient à $340\unite m$ pour un aller-retour.
\item $\Delta t=2d/v=6{,}8/340=0{,}020\unite s$. Vérifier l’unité de temps affichée.
\end{questions}
\par\smallskip\noindent\textbf{Pour comprendre.}\enspace Un écho de 2 s correspondrait à $340\times2/2=340\unite m$, pas 3,4 m. Pour 3,4 m : $\Delta t=2d/v=6{,}8/340=0{,}020\unite s$.\par
\end{corrige}
\end{minipage}
\par\noindent\begin{minipage}{\linewidth}
\phantomsection\label{cor-defi-19}\pcSolutionDefi{19}{Deux capteurs, deux milieux}\begin{questions}
\item A : $17\unite m$ ; B : $75\unite m$.
\item La vitesse de propagation dépend du milieu.
\item Le milieu, la vitesse utilisée et le caractère aller-retour du trajet, ainsi que les éventuels retards instrumentaux.
\end{questions}
\pcFigureEx{figures/cor-deux-milieux.pdf}{Même durée d’aller-retour, vitesses différentes : les deux distances ne sont pas égales.}
\end{minipage}
\par\noindent\begin{minipage}{\linewidth}
\phantomsection\label{cor-activite-19}\pcSousTitre{Pistes pour l’investigation}
Compter les motifs sur une durée connue : $f=N/\Delta t$. À durée égale, plus de motifs signifie une fréquence plus élevée et un son simple plus aigu. Changer l’amplitude ne change pas nécessairement la fréquence. Pour les distances, un écho correspond à deux trajets ; le tonnerre à un trajet simple dont le délai lumineux est négligé.\par\end{minipage}
